\chapterhead{48}{ORR-9}{Case Study: Cybersecurity and Information Security Risks}{1}{4}

\lohead{48}{a}{Provide examples of cyber threats and information security risks,
and describe frameworks and best practices for managing cyber risks.}

\orsub{Information security risk typology}

\begin{ordefbox}{Information security risk}
Information security risk encompasses both intentional and unintentional actions
that compromise the confidentiality, integrity, or availability of information.
This includes data leaks, hacking, and the loss or theft of both digital and
nondigital data.
\end{ordefbox}

\orsubb{Four-quadrant approach to risk categorisation}

\begin{center}
\begin{tikzpicture}[font=\fontsize{8.4}{10.5}\selectfont,
  q/.style={text width=5.4cm,align=center,rounded corners=3pt,inner sep=10pt,
            text=navy,draw=palenavy,line width=0.7pt},
  ax/.style={font=\sffamily\bfseries\fontsize{8.2}{10}\selectfont,text=navy}]

\node[ax] at (-3.2,0.95) {Intentional};
\node[ax] at ( 3.2,0.95) {Unintentional};
\node[ax,rotate=90] at (-6.6,0) {Internal};
\node[ax,rotate=90] at (-6.6,-2.5) {External};

\node[q,fill=palerust] (a) at (-3.2,0)    {Employee data theft};
\node[q,fill=paleblue] (b) at ( 3.2,0)    {Losing data};
\node[q,fill=palerust] (c) at (-3.2,-2.5) {External hacking};
\node[q,fill=paleblue] (d) at ( 3.2,-2.5) {Sending sensitive information externally};
\end{tikzpicture}
\end{center}
\orfigcap{Information security risks split by intent and by source.}

\orsub{Cybersecurity frameworks and standards}

\orsubb{1) NIST Cybersecurity Framework (NIST CSF)}

A voluntary framework to help businesses manage cyber risks. Five key steps:

\begin{center}
\begin{tikzpicture}[font=\fontsize{7.8}{9.6}\selectfont,
  st/.style={text width=2.55cm,align=center,rounded corners=3pt,inner sep=5pt,
             text=white,font=\sffamily\bfseries\fontsize{8.4}{10}\selectfont},
  sb/.style={text width=2.55cm,align=center,rounded corners=3pt,inner sep=6pt,
             fill=paleblue,draw=palenavy,line width=0.6pt,text=navy}]
\foreach \i/\x/\c/\t/\d in {%
  1/-6.5/navy/Identify/{Assets, IT equipment, software, data, policies},
  2/-3.25/teal/Protect/{Limit access, use security software, data backups},
  3/0/forest/Detect/{Monitor for unauthorised access, unusual activity},
  4/3.25/olive/Respond/{Develop a plan for cyberattacks, notify stakeholders},
  5/6.5/rust/Recover/{Repair impacted systems, restore data, improve procedures}}
{
  \node[st,fill=\c,anchor=north] (t\i) at (\x,0) {\t};
  \node[sb,anchor=north] (b\i) at (t\i.south) {\d};
}
\foreach \a/\b in {1/2,2/3,3/4,4/5}
  \draw[-{Stealth[length=4pt]},navy!35,line width=1.2pt]
    ($(t\a.east)+(0.02,0)$) -- ($(t\b.west)+(-0.02,0)$);
\end{tikzpicture}
\end{center}
\orfigcap{The five functions of the NIST Cybersecurity Framework.}

\orsubb{2) CIS Critical Security Controls}
\begin{itemize}
\item \textbf{Purpose:} provides detailed guidance for establishing cybersecurity
protocols.
\item \textbf{Features:} protects against widespread cyberattacks; complements
other standards such as NIST and ISO 27001.
\end{itemize}

\orsubb{3) ISO 27001}
\begin{itemize}
\item The key global standard for cybersecurity.
\item Provides guidance on information security risk management, operational
controls, and risk assessment.
\item Requires firms to develop an Information Security Management System (ISMS)
to manage risks.
\item Provides a checklist for firms to achieve certification in information
security management.
\end{itemize}

\orsub{Cybersecurity protection and monitoring}

\begin{orkeybox}
{\sffamily\bfseries\fontsize{8.6}{10}\selectfont\color{navy}Confidentiality,
Integrity, and Availability (CIA)}\par
\vspace{3pt}
The fundamental components of cybersecurity.
\begin{checks}
\item \textbf{Confidentiality} and \textbf{Integrity} relate to information
security (data protection).
\item \textbf{Availability} relates to business continuity (ensuring systems are
operational).
\end{checks}
\end{orkeybox}

\orsubb{Information controls}

\begin{lettered}
\item \textbf{Behavioural controls:} address human behaviour through training,
awareness programs, and password management.
\item \textbf{Technical controls:} address technical aspects of risk through
  \begin{itemize}
  \item \textbf{Preventative controls:} firewalls, encryption, patching.
  \item \textbf{Detective controls:} data loss prevention (DLP) solutions to
  identify unauthorised data transfers.
  \item \textbf{Mitigating controls:} backups to restore data in case of
  incidents.
  \end{itemize}
\item \textbf{Monitoring:} firms need to monitor cyber threats and breaches,
often through dedicated IT or information security departments.
\end{lettered}

% ---------------------------------------------------------------------
\lohead{48}{b}{Describe lessons learned from the Equifax case study.}

\orsub{Equifax case study}

\begin{orkeybox}
Equifax, a major credit bureau, suffered a huge data breach in 2017. Attackers
stole information on 147 million customers: names, addresses, Social Security
numbers, credit cards. The breach remained undetected for three months.
\end{orkeybox}

\orsubb{Equifax's five weaknesses}

\begin{lettered}
\item \textbf{Outdated cybersecurity practices.} Equifax relied on outdated
security frameworks and tools.
\item \textbf{Unpatched vulnerabilities.} An internal audit revealed a backlog of
unpatched vulnerabilities, including the exploited one.
\item \textbf{Patch management policy failure.} Equifax failed to enforce its
policies for timely security patches.
\item \textbf{Ineffective communication.} Communication about vulnerabilities
didn't reach relevant employees.
\item \textbf{Weak monitoring.} Expired SSL certificates hindered proper traffic
monitoring.
\end{lettered}

\begin{ornotebox}{Outcome}
Equifax faced loss of reputation and customer trust, and a \$700 million
settlement for fines and compensation. After this, Equifax overhauled its IT
security systems and management.
\end{ornotebox}
